\documentclass[11pt,a4paper]{article}

\usepackage[T1]{fontenc}
\usepackage{amsmath,amssymb,amsthm,mathtools}
\usepackage{booktabs}
\usepackage{enumitem}
\usepackage{geometry}
\usepackage{microtype}
\usepackage{listings}
\usepackage{hyperref}

\geometry{margin=25mm}
\hypersetup{
  colorlinks=true,
  linkcolor=blue,
  citecolor=blue,
  urlcolor=blue,
  pdftitle={The Maximum Determinant of a 55 by 55 Binary Circulant Matrix},
  pdfauthor={Qichao Wang and Daoyu Dong}
}
\urlstyle{same}
\setlength{\parindent}{0pt}
\setlength{\parskip}{5pt}
\setlength{\emergencystretch}{2em}
\lstset{
  basicstyle=\ttfamily\small,
  columns=fullflexible,
  keepspaces=true,
  breaklines=true,
  frame=single,
  framerule=0.3pt,
  xleftmargin=4pt,
  xrightmargin=4pt
}

\newtheorem{theorem}{Theorem}
\newtheorem{proposition}{Proposition}
\newtheorem{lemma}{Lemma}
\theoremstyle{remark}
\newtheorem*{remark}{Remark}
\DeclareMathOperator{\Res}{Res}
\DeclareMathOperator{\wt}{wt}
\newcommand{\Z}{\mathbb Z}
\newcommand{\F}{\mathbb F}
\newcommand{\D}{D_{01}(55)}
\newcommand{\M}{134694094094758395331307111329132}

\title{The Maximum Determinant of a $55\times55$ Binary Circulant Matrix}
\author{
  Qichao Wang\\
  Hebei University of Technology
  \and
  Daoyu Dong\\
  University of Electronic Science and Technology of China
  \thanks{Qichao Wang and Daoyu Dong contributed equally to this work.}
}
\date{}

\begin{document}
\maketitle

\begin{abstract}
For a binary word $a=(a_0,\ldots,a_{54})$, let $C(a)$ be the circulant
matrix with $(i,j)$ entry $a_{j-i}$, with subscripts modulo $55$.  We give a
computer-assisted exhaustive proof that
\[
  \D=\max_{a\in\{0,1\}^{55}}|\det C(a)|
  =\M.
\]
All maximizing words have weight $28$ and form one affine-equivalence class
under $i\mapsto ui+t\pmod {55}$, with $u$ a unit; the class contains
$2200$ words and has trivial stabilizer.  The proof reduces the $2^{55}$
words by complement duality, exact spectral-moment bounds, the
$5\times11$ cyclotomic decomposition, correlation-profile enumeration, and
exhaustive binary lifts.  The production certificate records $38,629,684$ ordered correlation
profiles, explicitly enumerates $6,845,230$ of them, and records $16,084$ lift
jobs containing $15,487,882,832$ joined words.  Exact Fourier reconstruction,
integer determinants, and cyclotomic resultants verify the displayed winner.
A clean-room audit independently matches the complete profile target set and
the lift task partition.  The existing six-plus-five lift run is a
split-dependent replay of the same native lift algorithm, not a second full
fiber-union implementation; that remaining completeness audit is reported
explicitly below.  Accordingly, this manuscript is a review draft and is not
labelled arXiv-ready until that audit is closed.  To the best of our
knowledge, the exact value at order $55$ had not previously been determined
in the sources checked for this work.
\end{abstract}

\section{Introduction}

Binary circulant matrices are a structured class of integer matrices whose
determinants can be evaluated spectrally, while the underlying words still
have a search space of size $2^n$.  Write
\[
  D_{01}(n)=\max_{a\in\{0,1\}^n}|\det C(a)|.
\]
The exact binary-circulant tables of Brent and Yedidia extend through order
$53$ in their updated work \cite{brent-yedidia}.  The order-$54$ computation
ALETHEIA-54 provides a subsequent exact computational reference
\cite{aletheia54}.  We determine the next order, $55=5\cdot11$.

The novelty statement is intentionally scoped.  The source and retrieval
record used for this manuscript reports no prior published exact
determination of $D_{01}(55)$ in the sources checked as of 9 September 2026;
the search is not an absolute-priority claim.  The checked records include
the OEIS sequence and b-file, the Brent--Yedidia sources, ALETHEIA-54, and
targeted searches for the exact integer and for order-$55$ binary circulants
\cite{oeis}.

The main result is the following.

\begin{theorem}[Global maximum and classification]
\label{thm:main}
For binary circulant matrices of order $55$,
\[
  D_{01}(55)=\M.
\]
Moreover, every maximizing word has weight $28$.  The maximizing words form
one orbit under the affine action
\[
  a_i\longmapsto a_{ui+t},
  \qquad u\in(\Z/55\Z)^\times,
  \quad t\in\Z/55\Z,
\]
and this orbit has $2200$ elements.
\end{theorem}

The canonical representative used throughout is
\[
\boxed{\texttt{0000000100011111011011001101011011010100011110101000111}}
\]
and has weight $28$.  It has stabilizer size one under the affine action.
Its complement is
\[
\texttt{1111111011100000100100110010100100101011100001010111000},
\]
of weight $27$, and its determinant is
\[
  129883590734231309783760428781663.
\]
The scope of Theorem~\ref{thm:main} is binary circulant matrices, not all
binary matrices and not unrestricted $\{\pm1\}$ matrices.

The proof architecture is worth making explicit:
\[
\begin{array}{c}
\text{all }2^{55}\text{ words}
\longrightarrow \text{complement and weight reduction}\\[2pt]
\longrightarrow \text{exact spectral bounds}
\longrightarrow \text{mod-5/mod-11 folds}\\[2pt]
\longrightarrow \text{correlation profiles}
\longrightarrow \text{binary lifts}\\[2pt]
\longrightarrow \text{exact determinant comparison}.
\end{array}
\]
Each arrow is a mathematical reduction with an exhaustive finite residue,
rather than a claim that an unstructured program happened to search far
enough.

\section{Binary circulants and Fourier formulation}

Let $a\in\{0,1\}^{55}$, let $k=\sum_{j=0}^{54}a_j$, and define
\[
  C(a)_{ij}=a_{j-i},
  \qquad
  f_a(x)=\sum_{j=0}^{54}a_jx^j.
\]
If $\zeta=\exp(2\pi i/55)$, the Fourier eigenvalues of $C(a)$ are
$f_a(\zeta^r)$, $0\leq r<55$, up to the harmless choice of frequency
orientation.  Since $a$ is real and $55$ is odd, the nonzero frequencies
pair by complex conjugation.  Thus, with
\[
  q_r=|f_a(\zeta^r)|^2,\qquad 1\leq r\leq27,
\]
we have
\begin{equation}
\label{eq:fourier-det}
  \det C(a)=k\prod_{r=1}^{27}q_r\geq0.
\end{equation}
Consequently the absolute value in the definition of $D_{01}(55)$ is
redundant for actual binary words, although absolute products are retained
when formal correlation profiles have not yet been shown realizable.

Define the periodic autocorrelation sequence
\[
  c_s=\sum_{t=0}^{54}a_ta_{t+s},
  \qquad s\in\Z/55\Z.
\]
Then $c_0=k$, $c_s=c_{55-s}$, and counting ordered pairs of distinct
support elements gives
\begin{equation}
\label{eq:corr-sum}
  2\sum_{s=1}^{27}c_s=k(k-1).
\end{equation}
The correlation sequence is the inverse Fourier transform of the squared
Fourier magnitudes.  Parseval therefore gives the two moment identities
\begin{align}
  2\sum_{r=1}^{27}q_r&=55k-k^2, \label{eq:first-moment}\\
  k^4+2\sum_{r=1}^{27}q_r^2
    &=55\left(k^2+2\sum_{s=1}^{27}c_s^2\right).
    \label{eq:fourth-moment}
\end{align}
The first identity and AM--GM immediately imply
\begin{equation}
\label{eq:amgm}
  \det C(a)\leq k\left(\frac{k(55-k)}{54}\right)^{27}.
\end{equation}

\section{Complement duality and affine symmetry}

Let $\bar a=1-a$.  At the zero frequency,
$f_{\bar a}(1)=55-k$.  At every nonzero $55$th root of unity,
\[
  f_{\bar a}(\zeta^r)=-f_a(\zeta^r),
  \qquad 1\leq r<55,
\]
because the sum of all $55$th roots in a nontrivial frequency is zero.
There are $54$ nonzero eigenvalues, so the signs cancel in the determinant.
For $0<k<55$ this proves
\begin{equation}
\label{eq:complement}
  k\det C(\bar a)=(55-k)\det C(a).
\end{equation}
The constant words are singular.  If
\[
  D_{\mathrm{circ}}(55,k)=
  \max_{\wt(a)=k}\det C(a),
\]
then complement is a bijection between weight strata and
\begin{equation}
\label{eq:weight-pair}
  D_{\mathrm{circ}}(55,55-k)
    =\frac{55-k}{k}D_{\mathrm{circ}}(55,k).
\end{equation}
It is therefore enough to handle $1\leq k\leq27$.  For the complementary
word, the first-moment bound takes the form
\begin{equation}
\label{eq:complement-amgm}
  \det C(\bar a)
    =(55-k)\prod_{r=1}^{27}q_r
    \leq(55-k)\left(\frac{k(55-k)}{54}\right)^{27}.
\end{equation}

Translations and multipliers preserve the determinant.  Translation changes
Fourier eigenvalues only by phases whose product is one, while multiplication
of indices by $u\in(\Z/55\Z)^\times$ permutes the nonzero frequencies.  The
affine group has order
\[
  55\varphi(55)=55\cdot40=2200.
\]
Reversal is the multiplier $u=-1$.  Orbit sizes must still be computed from
actual stabilizers; the group order alone does not imply that every orbit has
size $2200$.  Correlations are invariant under translation and are indexed
by the $27$ opposite-shift pairs.  On these pairs the effective unit action
has order $20$.

\section{Spectral-moment reduction}

The correlation identities turn the global problem into a finite collection
of integer moment problems.  Among $h$ nonnegative integers with fixed sum
$T$, write $T=hb+r$, $0\leq r<h$.  Their minimum square sum is
\[
  (h-r)b^2+r(b+1)^2.
\]
Indeed, transferring one unit from two entries that differ by at least two
strictly decreases the square sum.  Applying this to the $27$ independent
correlations, and combining it with the exact product bounds below, excludes
all lower weights $k\leq24$ at threshold $M$.

For clarity, the remaining necessary windows are listed here.  The entries
are for $\sum_{s=1}^{27}c_s^2$; the upper endpoint is the largest integer
that can still reach $M$ under the certified moment bound.
\begin{center}
\small
\begin{tabular}{@{}rrrr@{}}
\toprule
lower weight $k$ & complement weight & minimum square sum & maximum square sum\\
\midrule
25 & 30 & 3336 & 3338\\
26 & 29 & 3913 & 3919\\
27 & 28 & 4563 & 4570\\
\bottomrule
\end{tabular}
\end{center}
For $k=27$, the deviations $d_s=c_s-13$ have sum zero, so
$\sum d_s^2$ is even; the final row is equivalently a small deviation
window around the nearly balanced correlation $13$.

\subsection{Exact product bounds}

The fourth moment in \eqref{eq:fourth-moment} fixes the square sum of the
$q_r$.  The following elementary extremal calculation supplies a rigorous
upper bound on their product.  Consider $h$ nonnegative variables with sum
$S$ and square sum at least $L$.  Contracting variables toward their common
mean increases the product until the square-sum constraint is tight (or the
variance is zero).  At a positive constrained stationary point, Lagrange
multipliers give
\[
  \frac1x=\alpha+2\beta x,
\]
so there are at most two distinct values.  If the smaller value occurs $m$
times, put $\mu=S/h$ and $V=L-S^2/h$.  The two values are
\begin{equation}
\label{eq:two-values}
  \mu-\sqrt{\frac{V(h-m)}{hm}},
  \qquad
  \mu+\sqrt{\frac{Vm}{h(h-m)}}.
\end{equation}
Enumerating $m=1,\ldots,h-1$, together with the zero-variance and boundary
cases, exhausts the compact feasible set.  The implementation bounds the
square roots outward by integer arithmetic at a decimal scale of $10^{50}$.
No floating-point comparison decides a pruning step.

The certificate also imposes a uniform spectral cap.  Let
\[
  b=\left\lfloor\frac{k(k-1)}{54}\right\rfloor,
  \qquad \delta_s=c_s-b.
\]
For every nonzero frequency,
\begin{equation}
\label{eq:spectral-cap-formula}
  q_r=k-b+\sum_{s=1}^{27}\delta_s\,2\cos(2\pi rs/55).
\end{equation}
The rearrangement inequality pairs the sorted deviations with sorted cosine
coefficients.  There are only three coefficient multisets, represented by
$r=1,5,11$, according to the gcd class of the frequency.  Their rational
outward bounds give a cap $B$ valid for every $q_r$.

To maximize a product subject also to $0\leq x_i\leq B$, enumerate the
number of coordinates on the upper boundary.  The remaining interior
coordinates again have the form \eqref{eq:two-values}; infeasible sums,
negative roots, and roots above $B$ are discarded.  This is a finite
description of all extrema on the capped set.  The cosine intervals are
constructed from Machin's identity and rational Taylor remainders for
$\pi$ and cosine, and the stored dyadic intervals are independently checked
by the verifier.

\section{The \texorpdfstring{$5\times11$}{5x11} cyclotomic decomposition}

The factorization $55=5\cdot11$ splits the frequency space into the
cyclotomic sectors of orders $1,5,11,55$.  Their dimensions are
\[
  1,\quad4,\quad10,\quad40,
  \qquad 1+4+10+40=55.
\]
The nonzero sectors have respectively $2$, $5$, and $20$ conjugate pairs.
Let
\[
  N_m(a)=\Res(f_a,\Phi_m),
  \qquad m\in\{5,11,55\}.
\]
All three cyclotomic degrees are even, so the resultant convention used here
has positive sign, and conjugate pairing makes these norms nonnegative.  We
obtain
\begin{equation}
\label{eq:resultant-factorization}
  \det C(a)=kN_5(a)N_{11}(a)N_{55}(a).
\end{equation}

Use the Chinese remainder identification
\[
  \Z/55\Z\cong\Z/5\Z\times\Z/11\Z.
\]
For a word $a$, define the two fold vectors
\[
  r_i=\sum_{t\equiv i\pmod5}a_t,\qquad
  s_j=\sum_{t\equiv j\pmod{11}}a_t.
\]
Thus $0\leq r_i\leq11$ and $0\leq s_j\leq5$.  The order-$5$ norm
$N_5$ depends only on $r$, and $N_{11}$ only on $s$.  If
\[
  h_i^{(m)}=\sum_{t\equiv i\pmod m}c_t,
\]
then $h^{(m)}$ is the cyclic correlation of the corresponding fold.  The
sector moments are
\begin{align}
  S_1^{(m)}&=\frac{m\sum_i r_i^2-k^2}{2},\\
  S_2^{(m)}&=\frac{m\sum_i (h_i^{(m)})^2-k^4}{2}.
\end{align}

For fixed folds $r$ and $s$, the first moment in the primitive order-$55$
sector is
\begin{equation}
\label{eq:primitive-sector-sum}
  T=\frac{55k+k^2-5\sum_i r_i^2-11\sum_j s_j^2}{2}.
\end{equation}
Its norm is at most $(T/20)^{20}$.  A simpler necessary fold screen is also
useful.  If the selected small sector has $h=(m-1)/2$ pairs and sum $T_m$,
the remaining $27-h$ pairs have sum
\[
  R=\frac{k(55-k)}2-T_m.
\]
Any word whose complementary determinant is at least $M$ must satisfy
\begin{equation}
\label{eq:fold-screen}
  (55-k)N_m\left(\frac{R}{27-h}\right)^{27-h}\geq M.
\end{equation}
This is rational and exhaustive over the bounded integer fold vectors.

The retained fold representatives, up to their affine fold symmetries, are:
\begin{center}
\small
\begin{tabular}{@{}rrr@{}}
\toprule
$k$ & order-$5$ representatives & order-$11$ representatives\\
\midrule
25 & 17 & 321\\
26 & 27 & 560\\
27 & 26 & 660\\
\bottomrule
\end{tabular}
\end{center}
All unit orientations of the folded correlations are retained.  Independent
translations of the two fold vectors lift jointly by the Chinese remainder
theorem.  This reduction does not assert that every pair of fold profiles
has a binary lift; that remaining existence question is handled exhaustively
in Section~\ref{sec:lifts}.

\section{Correlation profiles and exact enumeration}

For each active weight $k\in\{25,26,27\}$, the exact sum in
\eqref{eq:corr-sum}, the square-sum window, the range $0\leq c_s\leq k$, and
the capped moment bound define a finite family of integer correlation
profiles.  The generator first visits sorted multisets and uses the
multinomial coefficient to recover the exact number of ordered profiles.
It then tests every coordinate assignment against both fold signatures and
the effective order-$20$ unit action.

The resulting certificate counts are shown below.  ``Allowed'' is the
recorded overapproximation of formal ordered profiles after the initial
integer screen.  ``Explicitly enumerated'' is completed work, not projected
throughput.  A formal profile is not treated as a word until it passes the
binary lift.
\begin{center}
\scriptsize
\setlength{\tabcolsep}{4pt}
\begin{tabular}{@{}lrrrr@{}}
\toprule
$k$ & allowed ordered profiles & explicitly enumerated & canonical profiles & lift targets\\
\midrule
25 & 407,277 & 2,925 & 55 & 55\\
26 & 32,178,654 & 816,102 & 22,375 & 281\\
27 & 6,043,753 & 6,026,203 & 155,558 & 1,980\\
\midrule
total & 38,629,684 & 6,845,230 & 177,988 & 2,316\\
\bottomrule
\end{tabular}
\end{center}

For a formal profile, the $q_r$ obtained from the discrete Fourier transform
need not all be nonnegative.  The profile screen therefore uses the safe
inequality
\begin{equation}
\label{eq:formal-profile-bound}
  \left|(55-k)\prod_{r=1}^{27}q_r\right|
  \leq (55-k)\left(\frac{\sum_r q_r^2}{27}\right)^{27/2},
\end{equation}
which is just Cauchy--Schwarz followed by AM--GM on the absolute values.
The signed product is reconstructed from two $61$-bit Fourier fields.  A
separate squared integer comparison proves that the product of the primes
exceeds the bound in \eqref{eq:formal-profile-bound}; hence no modular wrap
can turn a formal product into a false target.  Profiles whose exact absolute
product is below $M$ cannot contain an improving word.  Profiles at or above
the threshold are passed to the complete lift enumeration.

\section{Exhaustive binary lifts}
\label{sec:lifts}

Use the explicit torus coordinate
\[
  (i,j)\in\Z/5\Z\times\Z/11\Z
  \longmapsto 11i+45j\pmod{55}.
\]
The $11$ columns are five-bit masks, so there are $32$ possible column
states.  For every retained correlation target, all compatible oriented
pairs of fold vectors are enumerated after normalizing the independent row
and column translations.  This produces exactly $16,084$ lift jobs.

The production lift splits the eleven columns into five and six.  A partial
assignment records its word and a seven-coordinate signature: five row
counts and the two within-column correlation sums for shifts $11$ and $22$.
These statistics are additive across columns.  The left half is sorted by
its packed signature; each right half is joined to the unique complementary
signature.  The packing is injective over the bounded coordinate ranges, so
the join uses exact integer equality rather than a hash collision.  Every
joined word has the requested folds and the two stored correlations; the
remaining independent shifts are checked by rotated $55$-bit intersections
and population counts.

Nonnegative partial row counts and partial correlations above their targets
are safe pruning rules.  No floating-point operation is used in the lift.
The production run tests $15,487,882,832$ joined words and finds exactly
two normalized solutions.  The canonical verifier replays all jobs with a
six-plus-five split and compares joined-word counts and solution sets.  This
is a split-dependent cross-check because it reuses the same native lift
algorithm; it is not an independent full fiber-union proof.  A separate small
fiber is cross-checked against all $78,125$ direct column assignments using
three splits.  The clean-room audit therefore certifies the listed task
partition and the displayed witnesses, while leaving the independent union
certificate for every full order-$55$ fiber open.

\section{Winner and classification}

For the canonical winner $w$, the exact folded data and cyclotomic norms are
\begin{lstlisting}
weight = 28
mod-5 margin:  [5,4,8,6,5]
mod-11 margin: [3,1,4,2,3,1,1,3,3,3,4]
N5  = 131
N11 = 434039
N55 = 84603917386870862636041
\end{lstlisting}
Equation \eqref{eq:resultant-factorization} gives
$28\cdot131\cdot434039\cdot84603917386870862636041=M$.
The complement has the same three nontrivial norms and the zero-frequency
factor changes from $28$ to $27$, giving the determinant displayed after
Theorem~\ref{thm:main}.

The nonzero correlations of the winner have the distribution
\[
  13\text{ four times},\qquad
  14\text{ forty-six times},\qquad
  15\text{ four times}.
\]
Among the independent shifts $1\leq s\leq27$, the exceptional shifts are
$2,6$ with value $15$ and $7,15$ with value $13$.  For the complement, the
half-profile deviations from $13$ have square sum $4$.  If the word is
converted to signs by $x_t=2a_t-1$, its nonzero periodic autocorrelations
are $-5,-1,3$ with total squared energy $182$ over all $54$ nonzero shifts.

For completeness, the perfectly balanced weight-$27$ possibility would have
$c_s=13$ for every nonzero shift, corresponding to a cyclic
$(55,27,13)$ difference set.  It is impossible.  Indeed, with
$\alpha=f_a(\zeta_5)$, where $\zeta_5$ is a primitive fifth root, the ideal
correlations imply $\alpha\overline\alpha=14$.  The polynomial
$\Phi_5$ is irreducible modulo $2$, since $\operatorname{ord}_5(2)=4$;
hence $\mathbb Z[\zeta_5]/(2)$ is a field.  Reducing the norm equation
modulo $2$ forces $\alpha=2\beta$.  Then
$\beta\overline\beta=7/2$, impossible because a rational algebraic integer
is an integer.  This obstruction is also discussed in the design literature
\cite{braic}.  The actual winner is not being presented as a new difference
set or as a two-level almost difference set; the exact classification comes
from the exhaustive lift results.

The global classification now follows.  Any word with determinant at least
$M$ has a lower complementary weight $1\leq k\leq27$.  The exact weight
screen excludes $k\leq24$.  For $k=25,26,27$, its moment identities place
the correlation multiset in the enumerated finite family; the folded screens
place both fold vectors among the retained representatives; and a unit
transformation takes the paired correlations to one of the $2,316$ targets.
The production lift certificate contains a transformed copy of the word.
The two normalized lifted solutions both give complementary determinant $M$
and have the same affine canonical representative.  The implication from
the finite profile screen to all binary words is conditional on the full
fiber-union statement; the current review certificate does not silently
replace that statement by the split-dependent replay.

Finally, direct enumeration of the affine action on the canonical word gives
$2200$ distinct words.  Its stabilizer therefore has size one, and the
exhaustive lift classification shows that all equality cases lie in this one
orbit.  This proves Theorem~\ref{thm:main}.

\section{Computer-assisted certification}

The proof separates discovery from certification.

\begin{description}[leftmargin=!,labelwidth=31mm]
\item[Heuristic stage.] Random or local search was used only to find a strong
incumbent.  Search scores, partial jobs, and floating-point rankings are not
evidence of maximality.
\item[Rigorous stage.] Exact first- and fourth-moment bounds, rational
outward intervals, exhaustive folded-profile coverage, exhaustive binary
lifts, exact integer arithmetic, and deterministic certificate checks prove
the global statement.
\end{description}

The final winner is checked through three mathematical routes:

\begin{enumerate}[leftmargin=*,itemsep=2pt]
\item finite-field Fourier determinants reconstructed by signed CRT;
\item an exact integer matrix determinant using fraction-free elimination;
\item the product of the three exact cyclotomic resultants.
\end{enumerate}

The production verifier does not import the search or certificate-builder
module.  It re-derives the rational bounds, checks the Machin/Taylor cosine
enclosures, verifies retained fold norms by resultants, recompiles the
hash-identified native generators, regenerates the correlation profiles, and
checks task identities, multiplicities, exact output words, determinants, and
the affine classification.  Its optional full mode is a split-dependent
six-plus-five replay of the native lift algorithm.  The final global
certificate records the component SHA-256 hashes and the counts used in this
paper; the independent completeness audit is a separate, explicitly
qualified review layer.


\section{Completeness statements and review status}
\label{sec:completeness-status}

The global argument uses three logically distinct finite statements.  We
state them separately because reproducibility of a computation is not the
same as an independently checkable certificate of its coverage.

\begin{lemma}[Correlation-spectrum coverage]
For each active weight $k\in\{25,26,27\}$, the integer moment equations,
the fold screens, and the unit action reduce every necessary correlation
profile to one of the finite profile records used in the certificate.
\end{lemma}

For unambiguous machine-readable transcription, the corresponding
certificate integers are $2316$ correlation targets, $6845230$ explicit
profile evaluations, $16084$ lift tasks, and $15487882832$ joined words.

The clean-room profile evaluator
\path|audit/completeness/run_cleanroom_profile_audit_v4.py| reconstructs the
three ordered-profile families without importing the production verifier or
calling a production executable.  It reports exact agreement of the 16
profile-family set, exact agreement of the 2,316 target triples, and 982,589
canonical profile records below the incumbent.  This is the independently
checkable portion of the profile-spectrum reduction.

\begin{lemma}[Pruning soundness]
Every profile recorded below the incumbent in the exact profile-product
stage is rigorously pruned for the purpose of reaching $M$; every retained
profile is represented by a lift target and is not pruned by that exact
profile-product comparison.
\end{lemma}

The clean-room records make this statement inspectable as a machine-readable
table with fields \path|profile_id|, \path|profile_data|,
\path|bound_type|, \path|exact_upper_bound|, \path|incumbent|, and
\path|pruned|.  The separate rational boxed-moment rows remain part of the
production bound certificate and are checked by the canonical verifier.

\begin{lemma}[Lift completeness, conditional form]
If the binary lift of every one of the 2,316 target profiles exhausts its
entire margin fiber, then the preceding two lemmas imply that every binary
word with determinant at least $M$ occurs in the recorded lift output, and
the winner and affine classification in Theorem~\ref{thm:main} follow.
\end{lemma}

The condition in this last lemma is the unresolved item identified by the
external completeness review.  The independent lift checker
\path|audit/completeness/independent_lift_checker.py| proves that the 2,316
targets expand to exactly 16,084 distinct task descriptors, that all 16,084
task audit rows are present, and that the two displayed witnesses realize
their task correlations and determinants.  It intentionally reports
\path|independent_full_fiber_union = NOT_COMPLETED|: the existing five-plus-six
run and its six-plus-five replay share the native lift implementation.

As additional controls, clean-room brute-force regressions at orders $15$ and
$21$ pass complement normalization, factor folding, profile grouping, and
exact fiber reconstruction against the full domains.  A deterministic sample
of exactly $1,000,000$ order-$55$ words is classified with no unknowns, and an
independent affine audit gives orbit size $2200$, stabilizer size $1$, and the
two exact determinant values stated above.  These controls strengthen the
review record but do not discharge the conditional lift-completeness lemma.

\paragraph{Publication status.}
The exact theorem and its production certificate are retained without
weakening the claimed result.  The present source is nevertheless a review
draft, not an arXiv-ready submission, until an independent order-$55$ fiber
union certificate (or an equivalent branch-level partition/union proof) is
added and checked.

\section{Reproducibility and availability}

The production theorem is tied to commit
\texttt{e77a1aa6afea16fe59a35f9c0de93f385f9b8781}.  The repository contains
the exact certificate directory, the native source files, and the verifier.
From the repository root, the fast publication checks are:
\begin{lstlisting}
python scripts/verify_order55_global.py
python -m pytest
python -m pip check
\end{lstlisting}
The split-dependent six-plus-five lift replay is requested with
\begin{lstlisting}
python scripts/verify_order55_global.py --full-lifts
\end{lstlisting}
The native replay requires a C++20 compiler with unsigned \path|__int128| support.
The stored build commands and source hashes identify the exact native inputs;
host-specific paths and elapsed times may differ, but the discrete counts,
words, products, and classification must agree.

The source code and machine-readable certificates accompanying this
manuscript will be archived with the public version of this work.  A
permanent repository URL was not available at packaging time, so no
unresolved URL placeholder is included in the arXiv source.

Qichao Wang and Daoyu Dong contributed equally to this work.

\section{Conclusion}

We have determined the exact maximum determinant of a $55\times55$ binary
circulant matrix and classified every equality case.  The essential global
step is the exhaustive chain from weight reduction to folded correlation
profiles and then to binary torus lifts.  The factorization $55=5\cdot11$
provides the cyclotomic sectors that make this chain finite and auditable,
while the final exact arithmetic removes any dependence on numerical
rounding.  The resulting value is
\[
  D_{01}(55)=\M,
\]
with one affine class of $2200$ maximizers.

\appendix
\section{Previously certified small-weight strata}

The starting repository also contains an independent exact certificate for
$D_{\mathrm{circ}}(55,k)$ with $1\leq k\leq7$.  These values are not a
separate headline result; they are supporting fixed-weight strata and can be
paired with their complements using \eqref{eq:weight-pair}.
\begin{center}
\small
\begin{tabular}{@{}rr@{}}
\toprule
$k$ & $D_{\mathrm{circ}}(55,k)$\\
\midrule
1 & 1\\
2 & 2048\\
3 & 1564031349\\
4 & 6424318816256\\
5 & 12556202952818275\\
6 & 4342461217049330166\\
7 & 570999857161750276477\\
\bottomrule
\end{tabular}
\end{center}

\bibliographystyle{unsrt}
\bibliography{references}

\end{document}
